\begin{definition}[The three plaquettes along the two-step route]
    \label{def:peps_two_step_string_plaquettes}
    \lean{TNLean.PEPS.torusTwoStepStringPlaquette}
    \leanok
    Relative to the four-endpoint patch based at $v$, the successive plaquette
    bases are $v+(1,1)$, $v+(2,1)$, and $v+(2,2)$.
\end{definition}
\begin{theorem}[Literal transports and clockwise fluxes along the route]
    \label{thm:peps_two_step_string_route_holonomy}
    \lean{TNLean.PEPS.torusSecondStringCrossingOutput_up_transport,TNLean.PEPS.torusSecondStringCrossingOutput_right_transport,TNLean.PEPS.regularWalkHolonomy_torusTwoStepString_initial,TNLean.PEPS.regularWalkHolonomy_torusTwoStepString_intermediate,TNLean.PEPS.regularWalkHolonomy_torusTwoStepString_final,TNLean.PEPS.regularWalkHolonomy_torusTwoStepString_partners}
    \leanok
    For $w\ge8$, $H\ge7$, the three actual clockwise fluxes are
    \begin{align}
        (h,1,1) & \longrightarrow(1,h,1)
        \longrightarrow(1,1,ghg^{-1}).
    \end{align}
    In the second step all upward transports remain those of the intermediate
    assignment; only the rightward transport based at $v+(2,2)$ is changed,
    to $gh^{-1}g^{-1}$. All native sorting inversions at torus seams are
    included. At both intermediate and final stages, the two first-string
    endpoint loops and the distant second-string partner have respectively
    the unchanged clockwise fluxes $g^{-1},g,h^{-1}$.
    These are actual bond-assignment calculations auxiliary to
    \cite[Section~6.6]{Schuch2010PEPS}; they do not identify the two-step
    operation with the prescribed complete braid.
\end{theorem}
\begin{proof}
    \leanok
    \uses{def:peps_two_step_string_plaquettes,thm:peps_second_string_physical_crossing,thm:peps_initial_string_right_physical_step,thm:peps_torus_plaquette_holonomy,thm:peps_torus_swept_string_endpoint_holonomy}
    The second operation changes exactly its selected middle horizontal bond.
    The native directed tables therefore determine every factor of each
    clockwise plaquette loop. Multiplying those four factors in their actual
    order gives the displayed values and the three retained partner fluxes.
\end{proof}
